%=============================================================================!
% 16.10  THE VIBRATIONAL DENSITY OF STATES                                    !
%=============================================================================!
\section{The vibrational density of states}
\label{sec:ob-vdos}

Infrared and neutron spectra report, each weighting the atoms in its own
way, at which frequencies the atoms of a material move, and so does a
trajectory: an atom's velocity, taken apart
into frequencies, says how much of its motion happens at each. Summed
over the atoms this is the vibrational density of states, the spectrum
by which a learned model of motion is first compared with the simulation
it learned from\cite{Thiemann2026}.

\subsection*{The power spectrum of a signal}

For a signal $x(t)$ recorded over a long time $t_{\mathrm f}$, the power at the
frequency $\nu$ is $\lvert\hat x(\nu)\rvert^2/t_{\mathrm f}$, with $\hat x$ the
transform of \cref{eq:ob-ft} over the record. Writing the squared length
as the transform times its conjugate,
\begin{equation*}
   \lvert\hat x(\nu)\rvert^2 = \int_0^{t_{\mathrm f}}\!\!\int_0^{t_{\mathrm f}}x(t)\,x(t')\,
   \mathrm{e}^{-2\pi\mathrm{i}\nu(t - t')}\,\dd t\,\dd t' ,
\end{equation*}
and averaging over many such records turns $x(t)x(t')$ into the
autocorrelation function $C(t - t')$, which in a settled signal depends on
the separation alone. Counting the pairs of times by their separation
$s$, as for \cref{eq:ob-msd-vacf}, each $s$ occurs over a length $t_{\mathrm f} -
\lvert s\rvert$, so
\begin{equation*}
   \frac{\ens{\lvert\hat x(\nu)\rvert^2}}{t_{\mathrm f}} = \int_{-t_{\mathrm f}}^{t_{\mathrm f}}\Bigl(1 -
   \frac{\lvert s\rvert}{t_{\mathrm f}}\Bigr)C(s)\,\mathrm{e}^{-2\pi\mathrm{i}\nu s}\,
   \dd s \to \int_{-\infty}^{\infty}C(s)\,\mathrm{e}^{-2\pi\mathrm{i}\nu s}
   \,\dd s
\end{equation*}
once $t_{\mathrm f}$ is long against the time over which $C$ dies away. The power
spectrum is the Fourier transform of the autocorrelation function, the
theorem of Wiener and Khinchin. Since $C$ is even, the imaginary part of
$\mathrm{e}^{-2\pi\mathrm{i}\nu s}$ cancels between $s$ and $-s$, leaving
$2\int_0^\infty C(s)\cos(2\pi\nu s)\,\dd s$.

\subsection*{The density of states}

Applied to the mass-weighted velocity autocorrelation function per atom,
$C_m(t) = \ens{m\vb{v}(t_0)\cdot\vb{v}(t_0 + t)}$, the theorem gives the
vibrational density of states, scaled so that its area counts the three
directions of each atom:
\begin{equation}
   \mathcal{D}(\nu) = \frac{4}{\kB T}\int_0^\infty C_m(t)\cos(2\pi\nu t)\,
   \dd t , \qquad \int_0^\infty\mathcal{D}(\nu)\,\dd\nu = 3 .
   \label{eq:ob-vdos}
\end{equation}
With $\hat C_m(\nu) = 2\int_0^\infty C_m(t)\cos(2\pi\nu t)\,\dd t$ the
transform of $C_m$, even in $\nu$, $\mathcal{D} = 2\hat C_m/\kB T$, and the
area follows from the inverse transform at $t = 0$: $C_m(0) =
\int_{-\infty}^\infty\hat C_m\,\dd\nu = 2\int_0^\infty\hat C_m\,\dd\nu$, so
the area of $\mathcal{D}$ is $C_m(0)/\kB T$, and $C_m(0) = \ens{mv^2} =
3\kB T$ by equipartition. For a harmonic solid the motion is a sum of
normal modes (\cref{sec:os-modes}), each oscillating at its own
frequency $\nu_k$ with the energy $\kB T$ on average. Each mode adds
$\kB T\cos(2\pi\nu_kt)/N$ to $C_m$ per atom, a spike of area $1/N$ in
$\mathcal{D}$ at $\nu_k$, so $\mathcal{D}$ is the histogram of the mode
frequencies, $3N$ spikes for $N$ atoms making the area 3: the classical
density of states. In a liquid, atoms that wander contribute at zero
frequency, where \cref{eq:ob-vdos} and \cref{eq:ob-gk} give
\begin{equation*}
   \mathcal{D}(0) = \frac{4}{\kB T}\int_0^\infty C_m(t)\,\dd t =
   \frac{12\,mD}{\kB T} ,
\end{equation*}
so the density of states holds the diffusion coefficient too.

A run gives $C_m$ only out to some lag $t_{\mathrm c}$, and cutting the
integral off sharply there adds ripples to $\mathcal{D}$, as cutting
$g(r)$ at half the cell did to $S(G)$ in \cref{sec:ob-sq}.
\texttt{spectra.vdos} multiplies $C_m$ by the window $\cos^2(\pi t/
2t_{\mathrm c})$, which falls smoothly from 1 to 0, transforms it by the
discrete transform of a record extended to negative times, and scales the
result to the area 3. The window blurs each peak over about
$1/t_{\mathrm c}$; a longer $t_{\mathrm c}$ sharpens the peaks and needs a
longer run to keep the noise down.

\subsection*{A crystal, with and without friction}

The fcc crystal at \SI{21}{\kelvin}, \SI{50}{\pico\second} at fixed energy
with velocities every \SI{10}{\femto\second} and $C_m$ taken to
\SI{10}{\pico\second}, has peaks at 0.50, 0.75, 0.95, 1.15, 1.40 and
\SI{1.90}{\tera\hertz} (\cref{fig:ob-vdos}a). Its exact harmonic
frequencies come from the Hessian of the same 256-atom cell
(\cref{sec:os-modes}), by finite differences of the forces: three zero
eigenvalues, the translations of the whole cell, and 765 modes up to
\SI{1.991}{\tera\hertz}. A finite cell has a few distinct frequencies,
and their bare histogram has spikes that the run's windowed spectrum
cannot show, so the fair comparison passes the modes' own $C_m(t) \propto
\sum_k\cos(2\pi\nu_kt)$ through the same window and transform. The overlap
of the two spectra is 0.839 by the score defined below. They agree in
the positions of the peaks, but the run's tallest peak is 3.74 per
terahertz against the modes' 4.28, and the run puts 0.192 of its area of
3 above the highest harmonic frequency, against 0.011 for the modes
through the same window: the run's motion is not exactly harmonic. These
overlaps are taken below \SI{3}{\tera\hertz}, where the crystal's
spectrum lies.

Friction widens every peak, as Chapter~4 found for the response of a
driven spring. A mode
of angular frequency $\omega_0$ damped by Langevin friction $\gamma$ swings
at $\omega_{\mathrm d} = \sqrt{\omega_0^2 - \gamma^2/4}$ (\cref{sec:os-damped}),
and when the friction is small against $\omega_0$ its velocity correlation
decays as $\mathrm{e}^{-\gamma\lvert t\rvert/2}\cos\omega_{\mathrm d}t$, with a
term in $\sin\omega_{\mathrm d}t$ of relative size $\gamma/2\omega_{\mathrm d}$
left out; at the crystal's lowest peak, \SI{0.5}{\tera\hertz}, that size
is 0.32, so the form is rough there. Each half of the cosine,
$\frac12\mathrm{e}^{\pm\mathrm{i}\omega_{\mathrm d}t}$, transforms into
\begin{equation*}
   \frac12\int_{-\infty}^{\infty}\mathrm{e}^{-\gamma\lvert t\rvert/2}\,
   \mathrm{e}^{-\mathrm{i}\delta t}\,\dd t
   = \frac12\Bigl[\frac{1}{\gamma/2 + \mathrm{i}\delta}
   + \frac{1}{\gamma/2 - \mathrm{i}\delta}\Bigr]
   = \frac{\gamma/2}{\gamma^2/4 + \delta^2} ,
\end{equation*}
with $\delta = \omega\mp\omega_{\mathrm d}$ and $\omega = 2\pi\nu$, the
integral over positive times being $\int_0^\infty\mathrm{e}^{-(\gamma/2 +
\mathrm{i}\delta)t}\,\dd t$ and that over negative times its conjugate.
This is a peak at $\pm\omega_{\mathrm d}$, a Lorentzian, that falls to
half its height where $\delta = \pm\gamma/2$: a full width of $\gamma$ in
angular frequency, or $\gamma/2\pi$ in $\nu$. Under friction of
\SI{2}{\per\pico\second} that is \SI{0.318}{\tera\hertz}. Spreading the
fixed-energy spectrum by a Lorentzian of that width reproduces the
Langevin run's (\cref{fig:ob-vdos}b): the overlap rises from 0.797 for
the unspread spectrum to 0.950, and the tallest peak falls from 3.74 per
terahertz to 2.01 in the Langevin run and 2.13 in the spread
spectrum. A spectrum computed
under a thermostat that acts on each atom belongs to the thermostat as
much as to the material.

\subsection*{A liquid, and two kinds of atom}

The argon liquid's density of states (\cref{fig:ob-vdos}c) has no sharp
peaks, a broad band near \SI{0.5}{\tera\hertz}, and a value at zero of
1.659 per terahertz, against $12mD/\kB T = 1.651$ with the Green-Kubo $D$
of the same run. In the molten salt of \cref{sec:ob-charge} each kind of
ion has its own spectrum (\cref{fig:ob-salt}c), computed from that kind's
velocities alone: the cations' band peaks at 3.70 and the anions' at
\SI{3.50}{\tera\hertz}, and the anions, of 1.54 times the mass and nearly
the same $D$, have the larger value at zero, 0.160 per terahertz against
0.098. Spectra of this kind, one for each element, are how TrajCast's
predictions are compared with molecular dynamics\cite{Thiemann2026}.

The comparison is condensed into a number by the overlap score of
Schran and co-workers\cite{Schran2021},
\begin{equation*}
   \text{overlap} = 1 - \frac{\int\lvert\mathcal{D}_1 - \mathcal{D}_2\rvert
   \,\dd\nu}{\int\mathcal{D}_1\,\dd\nu + \int\mathcal{D}_2\,\dd\nu} ,
\end{equation*}
which is 1 for identical curves and 0 for curves that do not overlap at
all. The two ions of the salt score 0.790 against each other; the
crystal's run against its harmonic modes 0.839.

\subsection*{A stride's Nyquist limit}

A learned propagator predicts frames a stride $\Delta t$ apart, and its
spectrum can be computed only from those frames; so can that of the
reference simulation, if it is to be compared fairly. By
\cref{sec:ob-fourier}, frames every $\Delta t$ show frequencies only up
to $1/2\Delta t$. For the strides of TrajCast's tests, below
\SI{7}{\femto\second}, that is \SI{71.4}{\tera\hertz} or
\SI{2383}{\per\centi\metre}, about the \SI{2400}{\per\centi\metre} that the
paper quotes as the range it can resolve\cite{Thiemann2026}. A faster
vibration reappears folded below the limit.

A chain of four beads of \SI{14.5}{\amu}, each standing for a group
CH$_2$ or CH$_3$, shows this. Its bonds have the energy $\kappa_r(r -
r_{\mathrm{eq}})^2$ with $\kappa_r = \SI{5}{\electronvolt\per\angstrom\squared}$, its
angles likewise, and a torsion twists it about its middle bond;
\cref{sec:ob-fes} sets the model out in full. Twenty copies are run at
fixed energy after \SI{10}{\pico\second} at \SI{300}{\kelvin}, with each
copy's centre-of-mass velocity removed (\cref{fig:ob-stride}). From every
frame, \SI{1}{\femto\second} apart, the spectrum has peaks at 0.50, 4.00,
9.00 and \SI{12.25}{\tera\hertz} and two stretches of the bonds at 18.25
and \SI{21.00}{\tera\hertz}. The bond alone, a spring of stiffness
$2\kappa_r$, the slope of the force $-2\kappa_r(r - r_{\mathrm{eq}})$, between two
beads of reduced mass $m_{\mathrm r}$ (\cref{sec:os-bodies}), would vibrate
at $\sqrt{2\kappa_r/m_{\mathrm r}}/2\pi = \SI{18.36}{\tera\hertz}$. From
every 20th frame, with its limit at \SI{25}{\tera\hertz}, the spectrum is
the same, and the overlap with that from every frame, both cut at
\SI{25}{\tera\hertz}, is 1.000. From every 30th, sampled at
$1/\Dt = \SI{33.33}{\tera\hertz}$ with its limit at
\SI{16.67}{\tera\hertz}, the two stretches fold to $33.33 - 18.25 = 15.08$
and $33.33 - 21.00 = \SI{12.33}{\tera\hertz}$, and the spectrum has peaks
at 15.15 and \SI{12.37}{\tera\hertz}, where nothing moves. The overlap
falls from 1.000 at \SI{20}{\femto\second} to 0.899 at 25 and 0.795 to
0.797 from 30 to \SI{40}{\femto\second}, the change coming as the
stride passes $1/(2\times\SI{21.00}{\tera\hertz}) = \SI{23.8}{\femto
\second}$. Spectra from strided frames must therefore be compared with
reference spectra taken at the same stride, and peaks near or beyond the
limit read with suspicion.

\begin{figure}[htb]
\centering
\includegraphics{ch16_observables/vdos}
\caption{(a) The vibrational density of states of the fcc crystal at
\SI{21}{\kelvin} at fixed energy, against the histogram of the harmonic
frequencies of the same cell (bars) and those frequencies passed through
the same window and transform (dashed). (b) The crystal under Langevin
friction of \SI{2}{\per\pico\second}, against the fixed-energy spectrum
spread by a Lorentzian of full width $\gamma/2\pi$ (dashed). (c) Liquid
argon at fixed energy, with $12mD/\kB T$ at $\nu = 0$ (point).}
\label{fig:ob-vdos}
\end{figure}

\begin{figure}[htb]
\centering
\includegraphics{ch16_observables/stride}
\caption{The four-bead chain at fixed energy, 20 copies, velocities
every \SI{1}{\femto\second} with each copy's centre-of-mass velocity
removed. (a) The vibrational density of states from every frame
(dashed), from every 20th and from every 30th, each to its own Nyquist
frequency (dotted lines). (b) The overlap score of the spectrum from
every $\Delta t$-th frame with that from every frame, both cut at the
Nyquist frequency of $\Delta t$, against $\Delta t$; the dotted line
marks $1/2\nu_{\mathrm b}$, with $\nu_{\mathrm b}$ the frequency of the
highest peak.}
\label{fig:ob-stride}
\end{figure}

\subsection*{How often an atom tries}

Chapter~12 noted that how often an atom hops depends also on how often it
tries, and left the measurement to this chapter. A lithium atom in a
hollow of the model surface vibrates, and each swing towards a bridge is
an attempt. Its density of states (\cref{fig:ob-layers}c) peaks at 5.25
and \SI{5.00}{\tera\hertz} at 600 and \SI{1000}{\kelvin}, below the
\SI{5.87}{\tera\hertz} of small vibrations in the hollow, because the
hollow widens as the atom climbs. That frequency is $\nu_0 =
\sqrt{k/m}/2\pi$, with $k$ the curvature of the energy at the bottom of
the hollow (\cref{sec:os-small}). Near a hollow each of the model's three
ripples (\cref{sec:en-surface}) has $\cos(\vb{g}_k\cdot\vb{r}) \approx 1 -
(\vb{g}_k\cdot\vb{r})^2/2$, and the squared cosines of the angles between
$\vb{r}$ and the three $\vb{g}_k$, \ang{120} apart, add to $\frac32$, so the
energy rises as $\frac{3}{16}U_{\mathrm b}\lvert\vb{g}\rvert^2r^2$ and $k =
3U_{\mathrm b}\lvert\vb{g}\rvert^2/8 = 2\pi^2U_{\mathrm b}/a^2 =
\SI{0.9785}{\electronvolt\per\angstrom\squared}$. With six bridges out of each hollow,
the estimate $6\nu\,\mathrm{e}^{-\beta U_{\mathrm b}}$ with this $\nu$
gives 0.095 and 0.923 hops per picosecond, against 0.134 and 0.898
counted and the transition-state rates 0.162 and 1.257 of
\cref{eq:th-tst}. The estimate takes the attempt frequency from the
spectrum and ignores how wide the passage over each bridge is compared
with the hollow, and it lies within a factor of 1.41 of the counted
rates: the frequency of attempts and the Boltzmann factor of the barrier
give a rate's size, and the counted hops give its value.

\begin{keyidea}
The vibrational density of states is the Fourier transform of the
mass-weighted velocity autocorrelation function, with area three per
atom; for a harmonic solid it is the histogram of the normal-mode
frequencies, and at zero frequency it is $12mD/\kB T$. Friction widens
each peak by $\gamma/2\pi$. Frames a stride $\Delta t$ apart show
frequencies only up to $1/2\Delta t$ and fold faster ones below it, so
spectra are compared at the same stride, by an overlap score.
\end{keyidea}

\notebook{16}{set the length of the correlation window and the stride, and watch the spectrum}

\begin{selfcheck}
The spectrum of one atom of a liquid is computed with $C_m$ cut at
\SI{1}{\pico\second}. Roughly how close can two peaks be and still be
seen as two?
\end{selfcheck}
